\subsection{Set Partitioning with controllable locality}
\label{subsec:partition_benchmark}
To study a feasible space distinct from assignment, consider a collection
of tiles $S_j\subseteq\{1,\ldots,m\}$, one per binary variable $x_j$.
Set Partitioning selects disjoint tiles covering every element exactly once:
$A_{uj}=\mathbf1\{u\in S_j\}$ and $Ax=\mathbf1$.
Singleton tiles guarantee feasibility. We use $N=n^2$ variables; here $n$
is a size label, not the number of equalities or binary variables.
An exact counting sampler generates uniform feasible covers, without solving
the objective or constructing the TT; its recurrence is given in
\cref{app:partition_protocol}.

For transfer of the assignment-selected ADD rule, the frozen locality grid
has $m\in\{18,19,20\}$, $n\in\{13,14,15\}$ and
$w\in\{5,6,8,10,12,15,18\}$. Tiles have one to four elements and span at most
$w$ consecutive element positions. Each $(w,m,n)$ cell contains 15 instances:
945 in total, or 135 for each $w$. A fixed spatial variable order processes
tiles by their smallest element, with the corresponding singleton last.
It allows completed element constraints to leave the active frontier.
Small $w$ therefore limits the number of simultaneous charge states; increasing
$w$ broadens the exact-indicator rank profiles. Mean maximum ranks range from
24.0 at $w=5$ to 19,962.3 at $w=18$ over the balanced parameter grid
(\cref{tab:partition_ranks}). These are full-indicator ranks, not the ranks
of a TT learned from 500 points. Locality also changes the feasible covers,
so rank is a descriptive covariate rather than an isolated causal treatment.

The nonpolynomial objective is a linear cost plus entropic scenario risk:
\begin{equation}
\begin{aligned}
C_{\mathrm{ent}}(x)=\frac{c^\top x}{s}
 +\tau\log\left[\frac1{32}\sum_{a=1}^{32}
 \exp\left(\frac{W_a^\top x}{\tau s}\right)\right],\\
 s=\sqrt m.
\end{aligned}
\label{eq:partition_entropic}
\end{equation}
with $\tau=0.5$, independent $c_j\sim U[-1,1]$ and
$W_{aj}\sim\mathcal N(0,1)$. The second term smoothly emphasizes adverse
scenario costs; the factor $1/32$ makes it an empirical exponential-utility
risk rather than a sum depending on scenario count. Log-sum-exp is evaluated
stably. The transfer protocol fixes GDF at $w_{\mathrm{filter}}=0.4$, selected
on assignment, and compares it with Best Cost using matched initial and
external feasible draws, without retuning on this grid
(\cref{app:partition_protocol}). The rank-dependent benefit is a hypothesis
of this transfer experiment, not inferred from the SA comparisons below.

For the completed SA experiment we use a separate low-rank Set Partitioning
family: connected tiles on a $3\times n$ strip, with $n=10,12,14$ and
$m=3n$. These instances are not the locality grid. Their fixed spatial order
and short tiles yield mean maximum exact ranks 27.7, 36.1 and 45.4,
respectively. Both quadratic $x^\top Qx$ and the entropic objective are used.
The two constructions and their frozen data are distinguished in
\cref{app:partition_protocol,app:sa_protocol}.

